Storage is the only term in the groundwater flow equations that connects one time step to the next, so it is the term through which an adjoint solution propagates backward in time. This chapter describes that coupling and the selection between the specific-storage and specific-yield contributions.

\section*{Coupling between time steps}

The right-hand side at time step $k$ depends on the head at step $k-1$ through the storage terms. Differentiating the adjoint problem with respect to that earlier head gives the term that carries the later adjoint state backward, which is what an instantaneous measure drops (Chapter~\ref{ch:instantaneous}). Writing it per cell,

\begin{equation}
	\label{eqn:drhsdh}
	\frac{\partial b}{\partial h^{k-1}} = -\left(
	\frac{\partial Q_{SS}}{\partial h^{k-1}} +
	\frac{\partial Q_{SY}}{\partial h^{k-1}} \right),
\end{equation}

\noindent where $Q_{SS}$ and $Q_{SY}$ are the volumetric flow rates from specific storage and specific yield (L$^{3}$T$^{-1}$).

The quantity in equation~\ref{eqn:drhsdh} is formed while the forward model runs and stored with that time step. It is paired with the adjoint state of the \emph{following} step, so the saturation it requires is the one at the step where it is formed, which the following step sees as the previous saturation.

\section*{Selecting the contributions}

\mf{} does not apply both contributions at once. The selection is made on the saturation of the cell at the previous time step, and it depends on the options in effect \citep{langevin2017}.

\subsection*{Specific yield}

Specific yield enters the previous-head derivative only through the saturation. For a cell that is neither full nor dry, $S_F^{k-1} = (h^{k-1} - BOT)/\Delta z$ and so $\partial S_F^{k-1} / \partial h^{k-1} = 1/\Delta z$, giving

\begin{equation}
	\label{eqn:dsy}
	\frac{\partial Q_{SY}}{\partial h^{k-1}} = \frac{S_y A}{\Delta t^{k}},
\end{equation}

\noindent where $S_y$ is the specific yield (dimensionless), $A$ is the horizontal area of the cell (L$^{2}$), and $\Delta t^{k}$ is the length of the time step (T). Once the cell is full or dry the saturation stops changing with head, the derivative in equation~\ref{eqn:dsy} vanishes, and specific yield drops out.

\subsection*{Specific storage}

The specific-storage contribution depends on which formulation is active. Under the default formulation the previous-head derivative follows the previous saturation,

\begin{equation}
	\label{eqn:dss-default}
	\frac{\partial Q_{SS}}{\partial h^{k-1}} =
	\frac{SC1}{\Delta t^{k}} S_F^{k-1}, \qquad SC1 = S_s A \Delta z,
\end{equation}

\noindent where $S_s$ is the specific storage coefficient (L$^{-1}$), $SC1$ is the primary storage coefficient (L$^{2}$), and $\Delta z$ is the cell thickness (L). Under the \texttt{SS\_CONFINED\_ONLY} option, specific storage acts only where the cell is full, and the term is removed rather than scaled:

\begin{equation}
	\label{eqn:dss-confined}
	\frac{\partial Q_{SS}}{\partial h^{k-1}} = \begin{cases}
		\dfrac{SC1}{\Delta t^{k}}, & S_F^{k-1} = 1, \\[2ex]
		0, & S_F^{k-1} < 1.
	\end{cases}
\end{equation}

A cell that is not convertible is treated as full at all times, and equation~\ref{eqn:dss-default} reduces to $SC1/\Delta t^{k}$.

Equations~\ref{eqn:dsy} through \ref{eqn:dss-confined} are complementary on the same test. Specific storage acts where the cell was full, specific yield where it was not.

\section*{Sensitivity to the storage parameters}

The terms above carry the adjoint solution backward through time. Specific storage and specific yield are also parameters of the flow model in their own right, and a performance measure has a sensitivity to each. They are reported separately, because they describe different water: specific storage the water an aquifer yields by compressing while it stays full, specific yield the water it drains from the pore space as the water table falls. A model calibrated against one is not thereby calibrated against the other.

Both follow from the same flows. Differentiating the specific-storage flow of equation~\ref{eqn:drhsdh} with respect to $S_s$ gives

\begin{equation}
	\label{eqn:dqdss}
	\begin{split}
	\frac{\partial Q_{SS}}{\partial S_s} = \frac{A \Delta z}{\Delta t^{k}}
	\Big[ & S_F^{k-1} h^{k-1} - S_F^{k} h^{k} \\
	& + BOT \left( S_F^{k} - S_F^{k-1} \right) \\
	& + \frac{\Delta z}{2} \left( \left(S_F^{k}\right)^{2}
	- \left(S_F^{k-1}\right)^{2} \right) \Big],
	\end{split}
\end{equation}

\noindent and the specific-yield flow, which is the water held between the saturations at the two ends of the step, gives

\begin{equation}
	\label{eqn:dqdsy}
	\frac{\partial Q_{SY}}{\partial S_y} = \frac{A \Delta z}{\Delta t^{k}}
	\left( S_F^{k-1} - S_F^{k} \right).
\end{equation}

Both expressions need the saturation at each end of the step. Equation~\ref{eqn:dqdsy} is nothing but the difference between them, so a saturation taken twice from the same end reduces it to zero, and equation~\ref{eqn:dqdss} to the term in the heads alone.

Under \texttt{SS\_CONFINED\_ONLY} the saturation is one at or above the top of the cell and zero below it, and specific storage acts only where the cell is full, so equation~\ref{eqn:dqdss} reduces to

\begin{equation}
	\label{eqn:dqdss-confined}
	\begin{split}
	\frac{\partial Q_{SS}}{\partial S_s} = \frac{A \Delta z}{\Delta t^{k}}
	\Big[ & \left. \left( TOP - h^{k} \right) \right|_{S_F^{k} = 1} \\
	& + \left. \left( h^{k-1} - TOP \right) \right|_{S_F^{k-1} = 1} \Big].
	\end{split}
\end{equation}

The two halves are taken on the saturation at their own end of the step, so a cell that is full at the start and not at the end carries the first half and not the second.

\section*{Verification}

A seven by seven single-layer model was drawn down by a well over four time steps, leaving the cells between about half and four-fifths saturated. The sensitivity of the head at an observation cell to the pumping rate was compared with a finite-difference derivative taken by perturbing that rate.

\begin{table}[htb]
	\centering
	\caption{Error in the head sensitivity against a finite-difference
	derivative, before and after the contributions were selected as described
	in this chapter.}
	\label{tab:sto-verify}
	\small
	\begin{tabular}{lrr}
		\toprule
		Formulation &
		\shortstack[r]{Applying $SC1$\\throughout} &
		\shortstack[r]{Selecting on\\saturation} \\
		\midrule
		\texttt{SS\_CONFINED\_ONLY} & 10.7\% & 0.17\% \\
		Default                     & 2.6\%  & 0.15\% \\
		\bottomrule
	\end{tabular}
\end{table}

Applying $SC1$ at full cell thickness everywhere gives a partially saturated cell a storage term the forward model does not have, and Table~\ref{tab:sto-verify} shows the resulting error. The residual error after the correction is consistent with the nonlinearity of the finite-difference derivative itself.

The sensitivities to the storage parameters themselves were checked on the same model, against finite-difference derivatives taken by perturbing each parameter in turn.

\begin{table}[htb]
	\centering
	\caption{Sensitivity to each storage parameter as a fraction of a
	finite-difference derivative in that parameter, before and after the
	derivatives were taken over the saturations at both ends of the time step.}
	\label{tab:sto-parameters}
	\small
	\begin{tabular}{lrr}
		\toprule
		Parameter &
		\shortstack[r]{Saturation\\taken once} &
		\shortstack[r]{Saturation at\\both ends} \\
		\midrule
		Specific storage & 0.9844       & 1.0000 \\
		Specific yield   & not reported & 1.0010 \\
		\bottomrule
	\end{tabular}
\end{table}

Equation~\ref{eqn:dqdsy} is the difference between the saturations at the two ends of the step, so taking both from the same end leaves it identically zero and there is no specific-yield sensitivity to report. The saturation the first time step starts from is the one the initial heads imply, and reading it before the flow model has formed it, when the array still holds the value it was allocated with, is enough on its own to change the sign of the summed specific-storage sensitivity.

The error is largest under \texttt{SS\_CONFINED\_ONLY} because that option removes the term entirely rather than reducing it, so the difference between applying and not applying it is the whole term. The size of the error in a given model scales with $S_s \Delta z$ relative to $S_y$: where a model sets $S_s$ such that $S_s \Delta z$ is comparable to or larger than $S_y$, an unselected specific-storage term dominates the coupling between time steps.

\section*{Limitation}

The \texttt{ORIGINAL\_SPECIFIC\_STORAGE} option, which restores the formulation used before \mf{} version 6.1.3, is not supported. Its previous-head derivative carries a term in the previous head itself in addition to the saturation, and \adj{} reports the option as unsupported rather than approximating it.
